\section{From a branching block to identical boundary measurements}
\label{sec:nonuniqueness}

We now use the exact scalar block supplied by
Proposition~\ref{prop:branching-block}. Set \(\Omega=B(0,3)\).
Our objective is to construct a uniformly positive bounded scalar
conductivity \(\gamma\), equal to one near the boundary, and a nonzero
\(w\in H^1_0(\Omega)\cap L^\infty(\Omega)\), supported away from the
outer boundary, whose source satisfies
\[
 \int_\Omega\gamma\nabla w\cdot
       \nabla\bigl((u-u_*)\phi\bigr)=0
 \qquad\bigl(\phi\in C_c^\infty(\Omega)\bigr)
\]
for every bounded \(\gamma\)-harmonic function \(u\), with a constant
\(u_*\) depending on \(u\). For a sufficiently small nonzero real
\(\delta\), this identity will make the change from \(\gamma\) to
\((1+\delta w)^2\gamma\) preserve every boundary measurement.
We obtain it by nesting scalar blocks, forcing one common
average of \(u\) on all cell boundaries, and passing signed surface
charges to the limit.

All similarities below act on scalar conductivities by composition, so
their ellipticity bounds remain fixed at every depth.

\subsection{A repeatable geometry}

Set
\begin{equation}\label{nu:parameters}
 L=\frac85,\qquad h=\frac3{100},\qquad
 s=\frac{57}{100},\qquad a=\frac{81}{100}.
\end{equation}
The inequalities $2s>1$ and $2s^3<1$ will respectively ensure
convergence of the potential series and vanishing volume of the
limiting nested set. Define the rectangular annular prism
\begin{equation}\label{nu:parent}
 D=\left\{x\in\R^3:
 h<\max\bigl(\abs{x_1}/L,\abs{x_2}\bigr)<1,
 \quad \abs{x_3}<1\right\}.
\end{equation}
Let \(R(x_1,x_2,x_3)=(x_2,x_1,x_3)\), put
\[
 S_\pm x=sRx\pm a e_1,\qquad D_\pm=S_\pm D,
 \qquad B=D\setminus\bigl(\overline{D_-}\cup\overline{D_+}\bigr),
\]
and write \(\Sigma=\partial D\) and
\(\Sigma_\pm=\partial D_\pm\).

\begin{lemma}\label{nu:geometry}
The three boundary components of \(B\) are disjoint Lipschitz tori.
They have compatible piecewise smooth bi-Lipschitz product collars,
whose removal leaves a connected Lipschitz central region. Moreover,
\(\overline{D_-}\) and \(\overline{D_+}\) are disjoint compact subsets
of \(D\).
\end{lemma}

\begin{proof}
The outer boxes of the children are, respectively,
\[
 [-1.38,-0.24]\times[-0.912,0.912]\times[-0.57,0.57]
 \quad\hbox{and}\quad
 [0.24,1.38]\times[-0.912,0.912]\times[-0.57,0.57].
\]
They lie strictly inside the parent outer box. Both avoid the parent
axial hole, whose first-coordinate half-width is \(Lh=0.048\), and
the two boxes have separation at least \(0.48\).

For explicit collars, let \(q(\theta)=(q_1(\theta),q_2(\theta))\)
parametrize the unit square perimeter by rays, and let
\(\zeta(\psi)\) similarly parametrize the perimeter of
\([h,1]\times[-1,1]\) about its center
\(c_0=((1+h)/2,0)\). Both parameters have period \(2\pi\).
Writing
\[
 (r,z)=c_0+(1+\tau)\bigl(\zeta(\psi)-c_0\bigr),
\]
the map
\begin{equation}\label{nu:collar}
 (\theta,\psi,\tau)\longmapsto
 \bigl(Lr q_1(\theta),\,r q_2(\theta),\,z\bigr)
\end{equation}
is a product collar of \(\Sigma\) for sufficiently small
\(\abs{\tau}\). Indeed \(r\) stays strictly positive, and on
each pair of perimeter faces the map and its inverse have bounded
derivatives and nonzero Jacobian. The finitely many seams have measure
zero. The child collars are its images under \(S_-\) and \(S_+\).
Choose their widths small enough that they are disjoint.

For completeness, the central complement is connected. Above the child
boxes, for example at height \(x_3=0.8\), its horizontal section is the
connected parent rectangular annulus. A point between heights
\(-0.57\) and \(0.57\) in the complement of the children can move
vertically to that level: its horizontal position is outside the
child annular footprints or inside one of their axial holes. A point
below the child boxes first moves horizontally in the parent annulus
to a position outside both boxes and then moves upward. The same
argument applies after sufficiently thin collars are removed. Such
removal slightly shrinks the parent cross-sectional rectangle in
\((r,x_3)\) and slightly expands the two child rectangles; all the
containment margins above remain positive.

We also check the local boundary regularity for these particular
regions. Positive separation places each boundary point on just one
boundary of a rectangular annular prism. A face is locally a half-space,
and the convex edge and corner types are graphs of maxima or minima of
finitely many affine functions after choosing a diagonal direction.
For an inner top corner, the reentrant local model, after changes of
coordinates and signs, is
\(\{\max(x,y)>0,\ z<0\}\). On the line
\((x,y,z)=(x_0,y_0,z_0)+t(1,1,-1)\), membership is equivalent to
\[
 t>\max\bigl\{z_0,-\max(x_0,y_0)\bigr\}.
\]
The threshold is a Lipschitz function on the transverse plane
\(x_0+y_0-z_0=0\). Inner edges give the corresponding two-dimensional
max graph, and reversing the side of any of these boundaries reverses
the inequality without changing the graph. The remaining edge and
corner types give the same max/min description. Thus every boundary
point has a local Lipschitz graph. The small collar adjustment keeps this
rectangular form and the positive separation, so the central region
is Lipschitz as claimed.
\end{proof}

\begin{figure}[tb]
\centering
\begin{tikzpicture}[x=2.3cm,y=2.3cm,line width=0.45pt,
                    every node/.style={font=\small}]
 \filldraw[fill=black!6] (-1.6,-1) rectangle (1.6,1);
 \filldraw[fill=white] (-0.048,-0.03) rectangle (0.048,0.03);
 \filldraw[fill=blue!15] (-1.38,-0.912) rectangle (-0.24,0.912);
 \filldraw[fill=blue!15] (0.24,-0.912) rectangle (1.38,0.912);
 \filldraw[fill=black!6,line width=0.3pt]
   (-0.8271,-0.02736) rectangle (-0.7929,0.02736);
 \filldraw[fill=black!6,line width=0.3pt]
   (0.7929,-0.02736) rectangle (0.8271,0.02736);
 \node at (-0.81,0.49) {\(D_-\)};
 \node at (0.81,0.49) {\(D_+\)};
 \node at (0,0.65) {\(B\)};
 \draw[<->] (-1.6,1.15) -- (1.6,1.15)
   node[midway,above] {\(2L=3.2\)};
 \draw[<->] (-1.76,-1) -- (-1.76,1)
   node[midway,left] {\(2\)};
 \draw[-{Stealth[length=3pt]}] (-0.30,-1.25) -- (0,-0.025);
 \node[below,align=center] at (-0.30,-1.25)
   {parent axial hole};
 \draw[-{Stealth[length=3pt]}] (1.15,-1.25) -- (0.81,-0.02);
 \node[below,align=center] at (1.15,-1.25)
   {child axial hole};
\end{tikzpicture}
\caption{Top view of the cell geometry, at the exact planar scale.
The parent prism has \(\abs{x_3}<1\); the two shaded child prisms
have \(\abs{x_3}<s\). The small gray child holes belong to \(B\),
whereas the white central hole is outside \(D\). The available space
above and below the children connects the block.}
\label{nu:geometry-figure}
\end{figure}

Let \(\mu\) be the probability measure on \(\Sigma\) obtained
by pushing forward \((2\pi)^{-2}\dd\theta\,\dd\psi\) under
Equation~\eqref{nu:collar} at \(\tau=0\), and let
\(\mu_\pm=(S_\pm)_*\mu\). These measures have bounded positive
densities relative to surface measure on their respective fixed
boundaries. Apply Proposition~\ref{prop:branching-block} to \(B\)
with these three measures. We obtain a scalar conductivity
\(b\) and constants \(0<c_b\le C_b<\infty\) such that
\(c_b\le b\le C_b\) almost everywhere, with the exact property:
every \(b\)-harmonic field having separately constant boundary
voltages has conormal flux a constant multiple of the corresponding
probability measure on each component.

The ball \(\Omega\) strictly contains \(\overline D\), since the
parent outer box has maximal norm \(\sqrt{L^2+2}<3\).
For a finite word \(\nu\) over \(\{-,+\}\), let \(\abs\nu\)
denote its length, set \(S_\varnothing=\Id\), and define
\[
 S_{\nu\pm}=S_\nu\circ S_\pm,\qquad
 D_\nu=S_\nu D,\quad B_\nu=S_\nu B,\quad
 \Sigma_\nu=\partial D_\nu,\quad \mu_\nu=(S_\nu)_*\mu.
\]
Every \(S_\nu\) is a similarity of ratio \(s^{\abs\nu}\).
Its orthogonal part can reverse orientation, which does not affect
any scalar energy or change-of-variables formula. The sets \(B_\nu\)
are pairwise disjoint. Define
\begin{equation}\label{nu:gamma}
 \gamma(x)=
 \begin{cases}
  b(S_\nu^{-1}x),&x\in B_\nu\text{ for some }\nu,\\
  1,&\text{otherwise}.
 \end{cases}
\end{equation}
This is measurable and satisfies
\(c\le\gamma\le C\), where
\(c=\min(1,c_b)>0\) and \(C=\max(1,C_b)<\infty\).
The remaining nested set has volume zero, since the total volume of
the depth-\(j\) cells is
\begin{equation}\label{nu:null-set}
 2^j s^{3j}\abs D=(2s^3)^j\abs D\longrightarrow0,
 \qquad 2s^3=0.370386<1.
\end{equation}
The countably many boundary seams are also null sets. Their assigned
values in Equation~\eqref{nu:gamma} are immaterial.

\subsection{Surface charges and equal averages}

We first record both the exact source identity and its scaling.
For vectors in \(\R^3\), \(\abs p\) denotes Euclidean norm.

\Needspace{16\baselineskip}
\begin{lemma}\label{nu:source-lemma}
There is a constant \(K\), independent of \(\nu\), such that for
every \(p=(p_0,p_-,p_+)\) with \(p_0+p_-+p_+=0\), there exists
\(W_{\nu,p}\in H^1_0(\Omega)\cap L^\infty(\Omega)\), vanishing
off \(\overline{D_\nu}\) and constant on each child cell, with
\begin{equation}\label{nu:source-identity}
 \int_\Omega\gamma\nabla W_{\nu,p}\cdot\nabla\psi
 =p_0\int_{\Sigma_\nu}\operatorname{Tr}\psi\,\dd\mu_\nu
  +p_-\int_{\Sigma_{\nu-}}\operatorname{Tr}\psi\,\dd\mu_{\nu-}
  +p_+\int_{\Sigma_{\nu+}}\operatorname{Tr}\psi\,\dd\mu_{\nu+}
\end{equation}
for every \(\psi\in H^1(\Omega)\). It obeys
\begin{equation}\label{nu:source-bounds}
 \norm{W_{\nu,p}}_\infty\le K s^{-\abs\nu}\abs p,
 \qquad
 \int_\Omega\abs{\nabla W_{\nu,p}}^2
 \le K s^{-\abs\nu}\abs p^2.
\end{equation}
\end{lemma}

\begin{proof}
On the fixed block, let \(V=(V_0,V_-,V_+)\) be the constant voltage
vector and let \(p\) be the integrated outward conormal flux vector
of its harmonic extension \(v\). The map \(V\mapsto p\) is linear,
\(\sum_i p_i=0\), and
\[
 \int_B b\abs{\nabla v}^2=\sum_i V_i p_i.
\]
Thus its kernel consists exactly of common constants: zero flux implies
zero energy, and connectedness of \(B\) then makes \(v\) constant.
The induced map from voltage vectors modulo constants to zero-sum flux
vectors is an isomorphism of two-dimensional spaces. Normalize its
inverse by \(V_0=0\).

Extend \(v\) by \(V_\pm\) in \(D_\pm\) and by zero outside
\(\overline D\). Matching traces give an \(H^1\) function
\(W_p\) with compact support in \(\overline D\).
The exact block flux property gives Equation~\eqref{nu:source-identity}
at the root. This identity holds for every \(H^1\) test: equivalently,
one may first multiply the test by a smooth cutoff equal to one near
\(\overline D\). The boundary trace functionals are continuous,
since the measures have bounded surface densities. Invertibility of
the finite-dimensional map, the maximum principle for the bounded
boundary voltages, and the energy identity give
\[
 \norm{W_p}_\infty\le K\abs p,
 \qquad \int\abs{\nabla W_p}^2\le K\abs p^2.
\]

For \(j=\abs\nu\), set
\[
 W_{\nu,p}(x)=s^{-j}W_p(S_\nu^{-1}x),
\]
where \(W_p\) is understood to be zero outside \(\overline D\).
A similarity of ratio \(r\) multiplies energy and integrated flux
by \(r\) in dimension three when the scalar coefficient is copied
without changing its values. The extra voltage factor \(r^{-1}\),
here \(r=s^j\), therefore preserves the prescribed integrated
charges and gives the bounds in Equation~\eqref{nu:source-bounds}.
The coefficient inside the children is irrelevant because this
function has zero gradient there.
Finally,
\((S_\nu)_*\mu_\pm=\mu_{\nu\pm}\); hence the measure on each
shared boundary is exactly the same from its two adjacent blocks.
\end{proof}

\begin{lemma}\label{nu:equal-averages}
If \(u\in H^1(\Omega)\) is weakly \(\gamma\)-harmonic, then
\[
 m_\nu(u):=\int_{\Sigma_\nu}\operatorname{Tr}u\,\dd\mu_\nu
\]
has one common value \(u_*\) for all finite words \(\nu\).
If \(u\) is bounded, then \(\abs{u_*}\le\norm u_\infty\).
\end{lemma}

\begin{proof}
Test the weak equation for \(u\) with \(W_{\nu,p}\), and use
Equation~\eqref{nu:source-identity} with \(\psi=u\). This gives
\[
 p_0m_\nu(u)+p_-m_{\nu-}(u)+p_+m_{\nu+}(u)=0
 \qquad\text{whenever }p_0+p_-+p_+=0.
\]
The triple of averages is orthogonal to the zero-sum plane and is
therefore constant. Iterating proves the assertion. If \(u\) is
bounded, its traces have the same essential bounds: apply the trace
map to the vanishing positive-part truncations beyond those bounds.
Since every \(\mu_\nu\) is a probability measure, the last claim
follows.
\end{proof}

\subsection{A bounded potential carrying opposite limiting charges}

Give the two first-level cells charges \(q_-=1\) and \(q_+=-1\),
and recursively put \(q_{\nu-}=q_{\nu+}=q_\nu/2\). Thus
\begin{equation}\label{nu:charges}
 \abs{q_\nu}=2^{1-j}\quad(\abs\nu=j\ge1),
 \qquad \sum_{\abs\nu=j}\abs{q_\nu}=2.
\end{equation}
Define
\[
 Z_0=W_{\varnothing,(0,1,-1)},\qquad
 Z_j=\sum_{\abs\nu=j}
       W_{\nu,(-q_\nu,q_\nu/2,q_\nu/2)}\quad(j\ge1),
 \qquad
 w^{(N)}=\sum_{j=0}^{N-1}Z_j.
\]
The source identities telescope. For every integer \(N\ge1\) and
every \(\psi\in H^1(\Omega)\),
\begin{equation}\label{nu:finite-source}
 \int_\Omega\gamma\nabla w^{(N)}\cdot\nabla\psi
 =\sum_{\abs\nu=N}q_\nu
      \int_{\Sigma_\nu}\operatorname{Tr}\psi\,\dd\mu_\nu.
\end{equation}

\begin{lemma}\label{nu:potential}
The sequence \(w^{(N)}\) converges in both \(L^\infty(\Omega)\)
and \(H^1(\Omega)\) to a nonzero
\(w\in H^1_0(\Omega)\cap L^\infty(\Omega)\) that vanishes
almost everywhere off \(\overline D\).
\end{lemma}

\begin{proof}
At a fixed depth the cell interiors are disjoint. Combining
Equations~\eqref{nu:source-bounds} and~\eqref{nu:charges}, with a
larger constant \(K\) if necessary, gives
\begin{equation}\label{nu:level-bounds}
 \norm{Z_j}_\infty\le K(2s)^{-j},\qquad
 \norm{\nabla Z_j}_2^2\le K(2s)^{-j}.
\end{equation}
Since \(2s=1.14>1\), both
\(\sum_j\norm{Z_j}_\infty\) and
\(\sum_j\norm{\nabla Z_j}_2\) converge. The first controls the
\(L^2\) norms as well because \(\Omega\) is bounded. This proves
the asserted convergence. Every partial sum is in \(H^1_0(\Omega)\)
and vanishes off \(\overline D\); the same is true of its limit.

To show that \(w\ne0\), fix
\(\eta\in C_c^\infty(\Omega)\) equal to one on a neighborhood
of \(\overline{D_-}\) and zero on a neighborhood of
\(\overline{D_+}\). The positive separation of these compact sets
makes this possible. At every depth the signed total in the first
left cell is \(1\) and that in the first right cell is \(-1\).
Equation~\eqref{nu:finite-source} therefore gives
\[
 \int_\Omega\gamma\nabla w^{(N)}\cdot\nabla\eta=1.
\]
Strong \(H^1\) convergence yields
\begin{equation}\label{nu:nonzero}
 \int_\Omega\gamma\nabla w\cdot\nabla\eta=1,
\end{equation}
which proves nontriviality. No trace or pointwise value on the limiting
nested set has been used.
\end{proof}

\subsection{An exact transformation of all boundary measurements}

The transformation in this subsection is related to the conformal
energy identity used by Daud\'e, Kamran, and
Nicoleau~\cite[Proposition~1.2]{DaudeKamranNicoleau}.
Here Equation~\eqref{nu:nonzero} shows that $w$ has a nonzero source.
The decisive property is that this source vanishes on every product
$(u-u_*)\phi$ with $u$ bounded and harmonic and $\phi$ a smooth compactly
supported test. We prove that identity before transforming the
harmonic extensions.

Choose a real \(\delta\ne0\) such that
\(\abs\delta\norm w_\infty<1/2\), and set
\begin{equation}\label{nu:new-coefficient}
 \rho=1+\delta w,\qquad \widetilde\gamma=\rho^2\gamma.
\end{equation}
Then
\begin{equation}\label{nu:ellipticity}
 \tfrac12\le\rho\le\tfrac32,
 \qquad \tfrac c4\le\widetilde\gamma\le\tfrac{9C}4
 \quad\text{almost everywhere}.
\end{equation}
Both conductivities equal one in a neighborhood of \(\partial\Omega\).
They differ on a set of positive measure: by
Equation~\eqref{nu:nonzero}, \(w\) is not zero almost everywhere,
and positivity of \(\rho\) makes \(\rho^2=1\) equivalent to
\(w=0\).

\begin{lemma}\label{nu:weighted-identity}
Let \(u\) be a bounded weakly \(\gamma\)-harmonic function in
\(\Omega\), and let \(u_*\) be its common average from
Lemma~\ref{nu:equal-averages}. Then, for every
\(\phi\in C_c^\infty(\Omega)\),
\begin{equation}\label{nu:weighted}
 \int_\Omega\gamma\nabla\rho\cdot
       \nabla\bigl((u-u_*)\phi\bigr)=0.
\end{equation}
\end{lemma}

\begin{proof}
The product \((u-u_*)\phi\) belongs to \(H^1_0(\Omega)\),
so it is an admissible test in Equation~\eqref{nu:finite-source}.
Put \(M=\norm u_\infty\). For each depth-\(N\) cell choose
\(x_\nu\in\overline{D_\nu}\). The zero-average identity on its
surface gives
\[
 \int_{\Sigma_\nu}(\operatorname{Tr}u-u_*)\phi\,\dd\mu_\nu
 =\int_{\Sigma_\nu}(\operatorname{Tr}u-u_*)
       \bigl(\phi-\phi(x_\nu)\bigr)\,\dd\mu_\nu.
\]
Since \(\abs{\operatorname{Tr}u-u_*}\le2M\) almost everywhere on this surface,
the absolute value of the sum on the right-hand side of
Equation~\eqref{nu:finite-source} is at most
\begin{equation}\label{nu:diameter-error}
 4M\norm{\nabla\phi}_\infty\,\operatorname{diam}(D)\,s^N.
\end{equation}
Here we used Equation~\eqref{nu:charges} and
\(\operatorname{diam}(D_\nu)=s^N\operatorname{diam}(D)\).
This tends to zero. The fixed product test has square-integrable
gradient, so strong \(H^1\) convergence of \(w^{(N)}\), together
with boundedness of \(\gamma\), justifies passage to the limit on
the left. Multiplication by \(\delta\) proves
Equation~\eqref{nu:weighted}.
\end{proof}

\begin{proposition}\label{nu:identical-dn}
The full weak Dirichlet-to-Neumann operators of \(\gamma\) and
\(\widetilde\gamma\) on \(\Omega\) are equal.
\end{proposition}

\begin{proof}
First take \(f\in C^\infty(\partial\Omega)\), and let \(u\)
be its weak \(\gamma\)-harmonic extension. The weak maximum
principle, obtained by testing truncations beyond the boundary bounds,
gives \(\norm u_\infty\le\norm f_\infty\). Define
\begin{equation}\label{nu:transformed-solution}
 \widetilde u=u_*+\frac{u-u_*}{\rho}.
\end{equation}
The Sobolev chain rule applies to \(1/\rho\) on the interval
\([1/2,3/2]\), and the product rule applies because both factors
are bounded \(H^1\) functions. Hence \(\widetilde u\in H^1(\Omega)\),
and it equals \(u\) off \(\overline D\), in particular near
the outer boundary. Its flux is
\begin{equation}\label{nu:transformed-flux}
 \widetilde\gamma\nabla\widetilde u
 =\rho\gamma\nabla u-(u-u_*)\gamma\nabla\rho.
\end{equation}
This is an \(L^2\) vector field.

For \(\phi\in C_c^\infty(\Omega)\), use \(\rho\phi\in H^1_0(\Omega)\)
in the original weak equation to obtain
\[
 \int_\Omega\rho\gamma\nabla u\cdot\nabla\phi
 =-\int_\Omega\phi\gamma\nabla u\cdot\nabla\rho.
\]
Equation~\eqref{nu:weighted}, expanded by the product rule, also gives
\[
 \int_\Omega(u-u_*)\gamma\nabla\rho\cdot\nabla\phi
 =-\int_\Omega\phi\gamma\nabla\rho\cdot\nabla u.
\]
Subtracting proves that the divergence of the flux in
Equation~\eqref{nu:transformed-flux} is zero. All cross-gradient
products are integrable by Cauchy--Schwarz; boundedness of the other
factors also gives the stated \(L^2\) flux bound. Density of smooth
tests in \(H^1_0(\Omega)\) proves the full weak equation.
By uniqueness, \(\widetilde u\) is the
\(\widetilde\gamma\)-harmonic extension of \(f\).

Now fix arbitrary \(g\in H^{1/2}(\partial\Omega)\) and choose
an \(H^1\) extension \(V\) of \(g\). Multiply \(V\) by a smooth
cutoff that equals one near \(\partial\Omega\) and is supported
in an outer collar disjoint from \(\overline D\). The resulting
\(V_0\) still has trace \(g\), and its gradient is supported
where the two solution fluxes agree. Consequently the defining weak
pairings satisfy
\[
 \langle\Lambda_{\widetilde\gamma}f,g\rangle
 =\int_\Omega\widetilde\gamma\nabla\widetilde u\cdot\nabla V_0
 =\int_\Omega\gamma\nabla u\cdot\nabla V_0
 =\langle\Lambda_\gamma f,g\rangle.
\]
Both operators are bounded from \(H^{1/2}\) to \(H^{-1/2}\) by
the elliptic energy estimate. Since smooth boundary functions are
dense in \(H^{1/2}(\partial\Omega)\) on the ball, equality extends
from smooth \(f\) to all boundary traces. This last step does not
require extending the transformation in
Equation~\eqref{nu:transformed-solution} to unbounded solutions.
\end{proof}

\begin{proof}[Proof of Theorem~\ref{thm:main}]
Lemma~\ref{nu:geometry} permits the use of
Proposition~\ref{prop:branching-block}, and
Equation~\eqref{nu:gamma} constructs a uniformly positive bounded
scalar conductivity on the ball \(\Omega=B(0,3)\subset\R^3\).
Lemmas~\ref{nu:source-lemma}--\ref{nu:potential} produce the bounded
nonzero function used in Equation~\eqref{nu:new-coefficient}.
The second scalar conductivity has the positive bounds in
Equation~\eqref{nu:ellipticity}, agrees with the first near the
boundary, and differs from it on a set of positive measure.
Set $\gamma_0=\gamma$ and $\gamma_1=\widetilde\gamma$, with common
positive finite bounds obtained by taking the smaller lower bound
and the larger upper bound. Proposition~\ref{nu:identical-dn} proves
equality of their full weak Dirichlet-to-Neumann operators. Thus
these coefficients give the
claimed counterexample in dimension three, disproving the assertion
quantified over all \(n\ge3\).
\end{proof}

\begin{proof}[Proof of Corollary~\ref{cor:localized}]
Return to the notation of Corollary~\ref{cor:localized}: let $\Omega$
be its bounded connected Lipschitz domain, let $c,C$ be the common bounds
from Theorem~\ref{thm:main}, and let $B_1,\ldots,B_m$ be the prescribed balls.
Write $B_i=B(x_i,r_i)$ and set $U_i=B(x_i,r_i/2)$.
The closed balls $\overline U_i$ are separated from one another and
from $\partial\Omega$, even if some prescribed balls have tangent
boundaries. Let $a_0,a_1$ be the pair of Theorem~\ref{thm:main},
and put $\Phi_i(y)=x_i+s_i y$, where $s_i=r_i/6$.
This maps $B(0,3)$ onto $U_i$. Copy the scalar coefficients by
composition: $a_{i,j}=a_j\circ\Phi_i^{-1}$ on $U_i$.
If $q_{i,j}(h)$ is the minimum energy with trace
$h\in H^{1/2}(\partial U_i)$, change of variables gives
\[
 q_{i,j}(h)=s_i
 \ip{\Lambda_{a_j}(h\circ\Phi_i)}{h\circ\Phi_i}.
\]
Thus $q_{i,0}=q_{i,1}$. The coefficient values and their unit
boundary collars are preserved; no amplitude rescaling is needed.

Define $\gamma_\varepsilon=a_{i,\varepsilon_i}$ on $U_i$ and
$\gamma_\varepsilon=1$ elsewhere. The stated bounds follow because
the original unit collars imply $c\leq1\leq C$.
Let $E=\Omega\setminus\bigcup_i\overline U_i$, and denote by
$E_{\rm ext}(f,h_1,\ldots,h_m)$ the minimum unit-conductivity
energy on $E$ with outer trace $f$ and inner traces $h_i$.
The trace theorem and matching-trace $H^1$ gluing, together with
Equation~\eqref{eq:energy-minimum}, give
\[
 \ip{\Lambda_{\gamma_\varepsilon}f}{f}
 =\min_{h_i\in H^{1/2}(\partial U_i)}
 \left\{E_{\rm ext}(f,h_1,\ldots,h_m)
       +\sum_{i=1}^m q_{i,\varepsilon_i}(h_i)\right\}
 \qquad(f\in H^{1/2}(\partial\Omega)).
\]
Indeed, restriction of a global competitor gives one inequality,
and gluing the subdomain minimizers gives the other; a global
minimizer supplies an attaining tuple of interface traces.
The right side is independent of $\varepsilon$. Polarization
therefore proves equality of the full weak operators.

If two indices differ at position $i$, their conductivities differ
inside $U_i$ on a set of measure
$s_i^3\abs{\{a_0\ne a_1\}}>0$. This proves all $2^m$ choices
are distinct. Their common boundary response is not asserted to be
that of the constant conductivity one.
\end{proof}
